\documentclass[11pt]{article}
\usepackage[margin=1in]{geometry}
\usepackage[T1]{fontenc}
\usepackage{lmodern}
\usepackage{enumitem}
\usepackage[hidelinks]{hyperref}

\setlist[itemize]{leftmargin=1.5em,itemsep=0.25em,topsep=0.35em}
\setlist[enumerate]{leftmargin=1.8em,itemsep=0.5em,topsep=0.4em}
\setlength{\emergencystretch}{2em}
\newcommand{\code}[1]{\texttt{#1}}

\title{Modern AutoCrit to MetricSpace:\\Workflow and Pre--October 20 Plan}
\author{Jonathan Ma, Sijia Zhu, Jiatong Gao}
\date{}

\begin{document}
\maketitle

\section{Modern AutoCrit Workflow}

Modern AutoCrit converts free-text clinical-trial eligibility criteria into structured, reviewable records. A user supplies an NCT identifier, pastes eligibility text, or uploads a text-based protocol PDF. For an NCT identifier, the application retrieves the registered trial information from ClinicalTrials.gov. For a complete protocol, it locates the eligibility section and limits extraction to that material.

The system separates inclusion and exclusion criteria, breaks compound text into atomic criteria, and extracts clinically meaningful fields. These fields include the clinical domain, entity, comparator, value, unit, numerical interval, negation, temporal restriction, repetition, qualifier, logical relationship, and verbatim source text. Terminology normalization maps synonymous expressions and abbreviations to consistent concepts while preserving the original evidence.

The extraction result is displayed in a sortable human-review table. A reviewer can accept a row or reject and correct it. Reviewed criteria are then written to timestamped JSON and CSV library snapshots. The resulting records retain both normalized structure and source provenance, making them suitable for downstream comparison while keeping model-generated decisions subject to human verification.

In short, the current workflow is:

\begin{center}
Protocol or NCT record
$\longrightarrow$
atomic extraction
$\longrightarrow$
structured normalization
$\longrightarrow$
human review
$\longrightarrow$
versioned criterion library.
\end{center}

\section{What MetricSpace Builds on Top}

Modern AutoCrit answers the question, ``What eligibility criteria does this trial contain, and how can each criterion be represented computationally?'' MetricSpace begins with those reviewed outputs and asks, ``How similar, contradictory, overlapping, or restrictive are criteria across trials?''

This distinction matters because ordinary sentence embeddings often emphasize shared wording without representing the population defined by a criterion. For example, \texttt{age $\geq 18$} and \texttt{age $<18$} may have high cosine similarity even though they define disjoint age ranges. More complex criteria can differ through negation, exceptions, clinical severity, medical history, timing, or repeated measurements.

MetricSpace will use the normalized fields and source evidence from Modern AutoCrit to:

\begin{itemize}
    \item construct pairs of eligibility criteria from multiple protocols;
    \item compare a conventional embedding baseline with exact structured relationships where those relationships can be calculated;
    \item identify false-high and false-low embedding similarities;
    \item assign clinically interpretable pairwise relationship labels with LLM assistance and human validation;
    \item quantify uncertainty and preserve ambiguous cases rather than forcing a label; and
    \item ultimately use the reviewed labels to learn and evaluate a better dense representation of eligibility criteria.
\end{itemize}

The first stage is therefore dataset construction and baseline evaluation. Model training should begin only after the team has a sufficiently broad, reviewed set of criteria and a clear labeling guide.

\section{Using the Criteria Library Interface}

The Criteria Library page contains two controls intended to make team review and data transfer easier.

\subsection{Download Excel spreadsheet}

The \textbf{Download Excel spreadsheet} button creates a new \code{criteriaLibrary\_YYYYMMDD\_HHMMSS.xlsx} file from the library that is current at the moment of download. It does not modify the application or create another reviewed-library version. The export contains the immutable examples plus the latest cumulative set of human-reviewed trial criteria, even when the search box is filtering the table on screen.

The workbook has two worksheets:

\begin{itemize}
    \item \textbf{Criteria Library} contains one structured criterion per row. Its header is frozen, its columns can be filtered, and long source sentences are wrapped for review.
    \item \textbf{Field Guide} provides short definitions of the main fields. It is a quick reference, while the more complete labeling rules are given below.
\end{itemize}

The first column, \code{library\_source}, distinguishes the two kinds of rows:

\begin{itemize}
    \item \code{immutable\_base} identifies packaged examples. Their \code{trial\_id} is \code{EXAMPLE}. They are demonstrations and edge cases, not observations from an actual clinical trial and should not be counted as trial data in an analysis.
    \item \code{reviewed} identifies criteria that a person accepted or corrected in Modern AutoCrit. These are the rows to use when building the protocol-derived MetricSpace dataset.
\end{itemize}

The original five simple examples are retained, and the immutable set also contains harder examples involving opposite age thresholds, ``both'' versus ``either'' sex wording, investigator discretion, clinical exceptions, numerical intervals, repeated measurements, temporal windows, and compound OR statements. These examples are useful for teaching the schema and testing code, but they must remain separate from training, validation, and test results based on real trials.

Important workbook columns are grouped conceptually as follows:

\begin{itemize}
    \item \textbf{Identity and provenance:} \code{trial\_id}, \code{source\_id}, and \code{row\_id}. \code{row\_id} identifies an atomic structured record. Multiple rows can share a \code{source\_id} when one source sentence contains several atomic criteria.
    \item \textbf{Eligibility interpretation:} \code{criterion\_type}, \code{domain}, and \code{entity}. These state whether the record is inclusion or exclusion, what clinical class it belongs to, and which normalized concept it evaluates.
    \item \textbf{Quantitative or categorical constraint:} \code{comparator}, \code{value}, \code{unit}, and \code{value\_kind}. Categorical values generally use \code{present} or \code{absent}; quantitative rows retain the literal comparison.
    \item \textbf{Normalized numerical range:} \code{lower\_bound}, \code{lower\_inclusive}, \code{upper\_bound}, \code{upper\_inclusive}, and \code{interval}. Empty lower or upper bounds represent an unbounded side. The inclusive flags indicate closed endpoints.
    \item \textbf{Semantic modifiers:} \code{negated}, \code{temporal}, \code{repetition}, \code{qualifier}, \code{logical\_operator}, and \code{relations}. These must be considered when two criteria share an entity but impose different clinical meanings.
    \item \textbf{Evidence and audit:} \code{source} is the verbatim text against which the row should be checked; \code{reviewed\_at} records when a reviewed trial row was saved.
\end{itemize}

The \code{relations} cell joins derived semantic labels with a vertical bar for spreadsheet readability. For example, \code{HAS\_VALUE | HAS\_TEMPORAL | HAS\_MULTIPLIER} means that the criterion contains a comparison, a time restriction, and a repeated-measurement requirement. These labels summarize other columns; they are not substitutes for reading the underlying fields and source sentence.

For team labeling, save a working copy of the downloaded workbook and add new columns rather than replacing the structured columns. Suggested pair-labeling columns are \code{criterion\_a\_row\_id}, \code{criterion\_b\_row\_id}, \code{relationship}, \code{reviewer}, \code{rationale}, \code{uncertainty}, and \code{adjudicated\_label}. The application does not currently import edited Excel files, so the workbook is an export for review, pairing, and analysis rather than an alternate way to edit the governed library.

\subsection{Reset reviewed library}

The \textbf{Reset reviewed library} button returns the Criteria Library page to its immutable example state. After an explicit confirmation, it permanently removes generated \code{criteriaLibrary\_*} and historical \code{eavLibrary\_*} JSON/CSV snapshots from the governed runtime and configuration directories. It does not delete the immutable example file, extraction job records, downloaded spreadsheets, protocols, or unrelated files.

Reset should be used only when intentionally starting a new review dataset, clearing test extractions, or preparing a clean demonstration. Before resetting a dataset that may be needed, download its Excel workbook and copy its timestamped JSON/CSV snapshots to a clearly named archive. A downloaded workbook cannot restore the application library automatically.

\section{Manual Labeling Instructions}

Manual labeling has two separate layers. They must not be combined into one decision:

\begin{enumerate}
    \item \textbf{Criterion-level review:} verify that each individual criterion has the same structure and meaning as a reviewed Modern AutoCrit output row.
    \item \textbf{Pair-level labeling:} after both individual rows are correct, assign the relationship between criterion A and criterion B for MetricSpace.
\end{enumerate}

This distinction prevents an extraction mistake from being mislabeled as an embedding or clinical-similarity failure. If the team uses Chia criteria, the Chia annotations should first be mapped into the Modern AutoCrit representation described below. Chia supplies useful source annotations, but it does not supply the final MetricSpace pair label.

\subsection{Step 1: Create one atomic row}

Each row should describe one independently evaluable clinical restriction. Preserve the complete original sentence as source evidence, but split separately testable concepts into different rows. For example, ``GPL $>40$; MPL $>40$; APL $>50$'' should produce three rows because each laboratory measurement can independently determine eligibility. Do not split a phrase when its components only make sense together, such as ``blood pressure $>145/95$.''

When a source sentence produces multiple rows, keep the same source text or source identifier on those rows and record their Boolean relationship as AND or OR when the relationship is explicit. Do not silently infer a Boolean relationship when the source is ambiguous.

\subsection{Step 2: Verify the Modern AutoCrit fields}

For each atomic row, reviewers should check the following fields in order:

\begin{itemize}
    \item \textbf{Type:} choose \texttt{inclusion} or \texttt{exclusion} according to the protocol section and the role of the statement. Type is not the same as negation. An exclusion criterion can itself contain ``not'' or describe an exception.

    \item \textbf{Domain:} select the closest Chia-style domain used by Modern AutoCrit: \texttt{observation}, \texttt{condition}, \texttt{person}, \texttt{device}, \texttt{drug}, \texttt{visit}, \texttt{procedure}, or \texttt{measurement}. Use \texttt{measurement} for a measured quantity such as age, blood pressure, or a laboratory result; use \texttt{condition} for a diagnosis or clinical state. Apply the same convention consistently across the dataset.

    \item \textbf{Entity:} enter only the normalized clinical concept being evaluated. Remove the comparator, threshold, unit, negation, timing, frequency, and eligibility words. For example, the entity in ``age $\geq18$ years'' is \texttt{Age}, and the entity in ``no prior myocardial infarction'' is \texttt{myocardial infarction}.

    \item \textbf{Comparator:} record the literal relation \texttt{>}, \texttt{>=}, \texttt{<}, \texttt{<=}, or \texttt{=} for quantitative restrictions. Use \texttt{NA} when there is no quantitative comparison. Do not reverse the comparator because the criterion is an exclusion criterion.

    \item \textbf{Value:} preserve the threshold or categorical state. Quantitative examples include \texttt{40}, \texttt{145/95}, and \texttt{18}. For categorical criteria, use \texttt{present} when the concept is asserted and \texttt{absent} when its absence is required. Do not put explanatory prose into this field.

    \item \textbf{Unit:} retain the source unit, normalized consistently when possible. Use \texttt{NA} only when no unit applies or none is stated. Never invent a unit from clinical expectations.

    \item \textbf{Interval and endpoints:} for a numerical restriction, verify the lower bound, upper bound, and whether each endpoint is inclusive. ``Age 18 to 65 years'' becomes $[18,65]$ years; ``age $>18$'' becomes $(18,+\infty)$; and ``age $\leq65$'' becomes $(-\infty,65]$. The interval must agree with the comparator and must not be reversed for exclusion rows.

    \item \textbf{Negation:} mark negation only when the clinical entity itself is required to be absent. ``No history of myocardial infarction'' is negated and has value \texttt{absent}. In contrast, ``diabetes not due to steroids'' does not negate diabetes; the phrase describes an exception or qualifier.

    \item \textbf{Temporal:} separate any time window from the entity, such as ``within the last three months,'' ``at screening,'' or ``for at least six months.'' Retain the wording needed to interpret the restriction.

    \item \textbf{Repetition:} record multiplicity or frequency, such as ``on three occasions'' or ``two consecutive measurements.'' Do not merge this into the entity or numerical value.

    \item \textbf{Qualifier:} preserve a clinically meaningful severity, anatomical scope, exception, control status, or other modifier that is not represented elsewhere. Examples include ``clinically significant,'' ``in the study eye,'' and ``despite treatment.''

    \item \textbf{Logic:} mark \texttt{standalone}, \texttt{AND}, or \texttt{OR}. Logic describes how atomic rows derived from the same source combine; it does not describe whether the row is an inclusion or exclusion criterion.

    \item \textbf{Source:} keep the verbatim protocol or Chia text and a stable source identifier. Never rewrite the evidence to make it agree with the structured fields. A reader must be able to trace every row back to its source.
\end{itemize}

The displayed semantic relations such as \texttt{HAS\_VALUE}, \texttt{HAS\_NEGATION}, and \texttt{HAS\_TEMPORAL} are derived from these reviewed fields. Reviewers should correct the underlying comparator, negation, timing, repetition, qualifier, or logic rather than manually invent relation labels.

\subsection{Step 3: Accept, correct, or escalate}

Select the check mark only when the row is atomic, every structured field agrees with the source, and no material restriction has been lost. Select the cross when any field is wrong. Correct all affected fields, save the correction, and then treat the corrected row as accepted.

Do not guess when the source does not support a unique interpretation. Record the case as ambiguous and add a short note explaining what cannot be determined. Escalate clinically specialized terms, uncertain units, unclear logical scope, and statements whose inclusion/exclusion role conflicts with their wording. Reviewers should not use outside clinical knowledge to add facts that are absent from the source.

Useful reviewed examples are:

\begin{itemize}
    \item ``Inclusion: age 18 to 65 years'' $\rightarrow$ type \texttt{inclusion}, domain \texttt{person}, entity \texttt{Age}, interval $[18,65]$, unit \texttt{years}.
    \item ``Exclusion: GPL $>40$'' $\rightarrow$ type \texttt{exclusion}, domain \texttt{measurement}, entity \texttt{GPL}, comparator \texttt{>}, value \texttt{40}, interval $(40,+\infty)$. The comparator remains \texttt{>}.
    \item ``Exclude blood pressure $>145/95$ on three occasions'' $\rightarrow$ type \texttt{exclusion}, domain \texttt{measurement}, entity \texttt{blood pressure}, comparator \texttt{>}, value \texttt{145/95}, repetition \texttt{on three occasions}.
    \item ``No history of menopause'' $\rightarrow$ entity \texttt{menopause}, value \texttt{absent}, negation true; the type still comes from whether the protocol lists the statement as inclusion or exclusion.
\end{itemize}

\subsection{Step 4: Label the MetricSpace pair}

Only compare two criteria after both criterion-level records have been accepted or adjudicated. Review both the original text and the normalized fields. Assign exactly one relationship:

\begin{itemize}
    \item \textbf{Equivalent:} the criteria impose the same practical eligibility restriction despite wording differences.
    \item \textbf{A more restrictive than B} or \textbf{B more restrictive than A:} satisfying the narrower criterion implies satisfying the broader one, but not conversely. Always state the direction.
    \item \textbf{Partially overlapping:} some people could satisfy both criteria, while each criterion also admits people excluded by the other.
    \item \textbf{Contradictory or disjoint:} no participant can satisfy both restrictions under the stated interpretation.
    \item \textbf{Related but not directly comparable:} the criteria concern related clinical concepts but neither defines a clear subset, overlap, or contradiction using the available information.
    \item \textbf{Unrelated:} the criteria address different clinical concepts with no useful eligibility relationship.
    \item \textbf{Ambiguous or insufficient information:} the relationship cannot be determined reliably from the source and structured fields.
\end{itemize}

For every pair, retain criterion A and criterion B identifiers, the chosen relationship, reviewer identifier, a brief evidence-based rationale, and an uncertainty or escalation note when applicable. Numerical interval relationships should be calculated from the endpoints before assigning the label. For complex clinical sentences, an LLM may propose a relationship and rationale, but the human reviewer remains responsible for the accepted label.

If multiple reviewers label the same pair, they should work independently before discussing disagreements. Preserve both original labels and record the final adjudicated label; do not overwrite disagreement, because reviewer agreement and ambiguity are important evaluation results.

\section{Plan Before October 20}

\begin{enumerate}
    \item \textbf{Select two or three therapeutic areas.}
    \begin{itemize}
        \item Prefer areas with different criterion patterns, such as oncology, cardiovascular disease, and metabolic disease.
        \item Keep enough trials within each area to support comparisons within the same therapeutic area.
    \end{itemize}

    \item \textbf{Build a trial cohort.}
    \begin{itemize}
        \item Begin with approximately 10--15 trials per therapeutic area, or 30--40 trials in total.
        \item Record each NCT identifier, indication, phase, intervention, protocol availability, and reason for inclusion.
        \item Do not select only trials with unusually simple eligibility criteria.
    \end{itemize}

    \item \textbf{Run Modern AutoCrit.}
    \begin{itemize}
        \item Extract and human-review the eligibility criteria for the selected trials.
        \item Save the structured outputs and their source evidence.
        \item Record extraction failures, reviewer corrections, ambiguous criteria, and terminology-normalization problems.
        \item Freeze a reviewed dataset version for subsequent experiments so the baseline is reproducible.
    \end{itemize}

    \item \textbf{Construct criterion pairs.}

    The initial pair set should contain:
    \begin{itemize}
        \item clearly equivalent criteria;
        \item different wording with equivalent meaning;
        \item numerical overlap, subset, and disjoint cases;
        \item negation and exception cases;
        \item temporal and repetition differences;
        \item clinically related but operationally different criteria; and
        \item unrelated controls.
    \end{itemize}

    \item \textbf{Establish an embedding baseline.}
    \begin{itemize}
        \item Embed the raw criterion text with one standard clinical or general sentence-embedding model.
        \item Calculate cosine similarity for the selected criterion pairs and nearest neighbors.
        \item Identify false-high-similarity and false-low-similarity cases.
        \item Include exact examples such as \texttt{age $\geq18$} versus \texttt{age $<18$}, but also evaluate complex sentences collected from the protocols.
    \end{itemize}

    \item \textbf{Begin human/LLM labeling.}

    Define the labeling guide before using pair judgments for model training. Initial labels should include:
    \begin{itemize}
        \item equivalent;
        \item one criterion is more restrictive;
        \item partially overlapping;
        \item contradictory or disjoint;
        \item related but not directly comparable;
        \item unrelated; and
        \item ambiguous or insufficient information.
    \end{itemize}

    The LLM may propose a label and rationale, but a human reviewer should accept, correct, or mark the judgment as unresolved. Directional labels must identify which criterion is more restrictive.
\end{enumerate}

\section{Expected October 20 Progress}

The project does not need a completed learned metric by October 20. A strong progress report should contain:

\begin{itemize}
    \item a documented therapeutic-area and trial-selection strategy;
    \item a cohort table containing approximately 30--40 candidate trials;
    \item an initial set of reviewed Modern AutoCrit outputs with frozen version information;
    \item a diverse set of criterion pairs and a draft human-labeling guide;
    \item cosine-similarity baseline results from one embedding model;
    \item several clear examples showing where the baseline succeeds and fails; and
    \item a description of how LLM judgments, human corrections, and uncertainty will later become training data for MetricSpace.
\end{itemize}

These results establish the need for MetricSpace and provide the data foundation for the agentic labeling workflow and later representation-learning experiments.

\end{document}
